\documentclass[11pt]{article}
\usepackage[T1]{fontenc}
\usepackage[utf8]{inputenc}
\usepackage{lmodern,amsmath,amssymb,amsthm,microtype}
\usepackage[margin=1in]{geometry}
\usepackage[colorlinks=true,linkcolor=blue,citecolor=blue,urlcolor=blue]{hyperref}
\hypersetup{pdftitle={Polynomial decay for the Korteweg-de Vries equation at length 14 pi},pdfauthor={Kaiwen Jin}}
\newtheorem{theorem}{Theorem}[section]
\newtheorem{lemma}[theorem]{Lemma}
\newtheorem{proposition}[theorem]{Proposition}
\theoremstyle{remark}\newtheorem{remark}[theorem]{Remark}
\numberwithin{equation}{section}
\title{Polynomial decay for the Korteweg--de Vries equation at length $14\pi$}
\author{Kaiwen Jin\\\small\href{mailto:yui.ui.siki@gmail.com}{yui.ui.siki@gmail.com}}
\date{October 2026}
\begin{document}
\maketitle
\begin{abstract}
We prove uniform small-data polynomial decay for the real Korteweg--de Vries
equation on an interval of length $14\pi$, with homogeneous Dirichlet boundary
conditions and zero Neumann datum at the right endpoint. This critical length
has three undamped linear modes, including a zero frequency. Every sufficiently
small real $L^2$ energy solution satisfies $\|y(t)\|_2\le C(1+t)^{-1/2}$,
with constants independent of the datum. The quadratic response is made
periodic only on the exponentially stable subspace. Resonant terms in the
undamped subspace may grow, but their boundary observation vanishes. An explicit
nonzero double-frequency observation then gives a uniform quartic energy loss,
which yields the all-time decay estimate.
\end{abstract}

\section{Introduction}\label{sec:intro}
The Korteweg--de Vries equation on a bounded interval is a basic example in
the boundary control of dispersive equations. Homogeneous boundary conditions
can dissipate its $L^2$ energy, even though the equation has no interior damping.
The obstruction to linear stabilization is particularly concrete: at certain
interval lengths a finite-dimensional space of linear solutions has zero
boundary loss. These critical lengths, identified in Rosier's controllability
analysis \cite{Rosier}, are
\[
 \mathcal N=\left\{2\pi\sqrt{\frac{k^2+k\ell+\ell^2}{3}}:
 k,\ell\in\mathbb N^*\right\}.
\]
Small nonlinear solutions are globally well posed in the usual energy class;
we use the realization of Coron--Cr\'epeau \cite{CC} and the global statement
in \cite[Lemma 3.2]{NX25}. The issue here is a quantitative decay rate for
that flow when linear boundary loss misses some directions.

Nguyen \cite[Theorem 1.1]{Nguyen} obtains decay at least of order $t^{-1/2}$
when the undamped space has dimension one, and for certain even-dimensional
undamped spaces satisfying a nondegeneracy condition. His argument expands
the flow over a fixed time interval and detects a quadratic boundary response.
A zero critical frequency interacting with a nonzero one can prevent the full
quadratically forced solution from having a bounded separated form
\cite[Remark 2.6]{Nguyen}. The case of odd undamped dimension greater than one lies outside those decay theorems. Polynomial decay in that case is explicitly
left open in \cite[Section 2.4.5]{NX25}. Recent estimates near critical lengths
\cite[Theorem 1.6]{NX26} exclude the exact critical value in that theorem and
do not settle the present question.

We consider the concrete length $L=14\pi$. Its critical pairs are
$(7,7),(2,11),(11,2)$, giving one stationary mode and a conjugate oscillatory
pair. For the real equation
\begin{equation}\label{eq:kdv}
 \begin{cases}
 y_t+y_x+y_{xxx}+yy_x=0,&t>0,\quad 0<x<L,\\
 y(t,0)=y(t,L)=y_x(t,L)=0,&t>0,\\
 y(0,x)=y_0(x),&0<x<L,
 \end{cases}
\end{equation}
we prove the following uniform bound. An energy solution means the usual
mild solution in
$C([0,\infty);L^2(0,L))\cap L^2_{\mathrm{loc}}([0,\infty);H^1_0(0,L))$;
Neumann boundary values are interpreted through the standard $L^2$-in-time
hidden trace, recalled in Section~\ref{sec:linear}.

\begin{theorem}\label{thm:main}
Let $L=14\pi$. There exist $\delta,C>0$ such that every real energy solution
of \eqref{eq:kdv} with $\|y_0\|_{L^2(0,L)}\le\delta$ satisfies
\begin{equation}\label{eq:main}
 \|y(t)\|_{L^2(0,L)}\le C(1+t)^{-1/2}\qquad\text{for every }t\ge0.
\end{equation}
The constants are chosen independently of the initial datum.
\end{theorem}

This supplies polynomial decay at an exact critical length with a
three-dimensional undamped space. The estimate allows every sufficiently
small real $L^2$ datum, including arbitrary mixtures of the stationary,
oscillatory and stable components. In particular it gives an affirmative
answer at $14\pi$ to the odd-dimensional decay question above. We do not
claim the same conclusion for all critical lengths or an optimal rate.

The mixed resonance is dealt with by retaining the entire center forcing
through Duhamel's formula. Only the stable part of the quadratic response is
chosen periodic. The remaining center motion can have secular growth, but
its boundary derivative is identically zero. The observation is therefore a
finite temporal Fourier sum. Its constant term detects the stationary mode;
an explicit double-frequency coefficient detects the oscillatory pair.
Orthogonality over one period gives a positive lower bound uniformly over
all real center directions. A fixed late window removes the stable
homogeneous transient. The quadratic expansion has a cubic remainder in the
boundary observation norm for rough $L^2$ data, giving a quartic energy loss
on a single, datum-independent interval. Iterating that loss proves
\eqref{eq:main}. The fixed-window comparison follows Nguyen's approach
\cite[Section 5]{Nguyen}; allowing the invisible resonant response is the
change needed here.

Section~\ref{sec:linear} determines the undamped subspace and proves
exponential decay on its orthogonal complement. Section~\ref{sec:quadratic}
constructs the quadratic observation and proves its uniform lower bound.
Section~\ref{sec:nonlinear} establishes the nonlinear energy loss and proves
the theorem. All spatial norms below are on $(0,L)$; complex notation is used
only for linear calculations, and $\langle f,g\rangle=\int f\overline g$.
\section{Linear evolution and its undamped subspace}\label{sec:linear}
Put $L=14\pi$, and write $H=L^2(0,L;\mathbb C)$ with inner product
$\langle f,g\rangle=\int_0^L f\overline g\,dx$, linear in its first argument.
The real subspace will be used for the nonlinear equation. Let
\begin{equation}\label{eq:generator}
 Av=-v'-v''',\qquad D(A)=\{v\in H^3(0,L):v(0)=v(L)=v'(L)=0\}.
\end{equation}
Rosier's Proposition 3.1 \cite{Rosier} gives a contraction semigroup $S(t)$
on the real $L^2$ space. Extending it by real and imaginary parts gives the
complex contraction semigroup with precisely the domain in \eqref{eq:generator}.
On $D(A)$ define $\mathcal Bv=v'(0)$. Integration by parts gives
\begin{equation}\label{eq:dissipative}
 \operatorname{Re}\langle Av,v\rangle=-\tfrac12|\mathcal Bv|^2.
\end{equation}
The adjoint is $A^*v=v'+v'''$ on $\{v\in H^3:v(0)=v(L)=v'(0)=0\}$.
The semigroup difference quotient shows that $S(t)$ preserves $D(A)$,
commutes there with $A$, and evolves an eigenvector $Av=\lambda v$ as
$S(t)v=e^{\lambda t}v$.

\subsection{Mild solutions and boundary estimates}
For a fixed $T>0$ set
\[
 X_T=C([0,T];H)\cap L^2(0,T;H^1_0(0,L)),\qquad
 \|v\|_{X_T}=\|v\|_{C_tH}+\|v\|_{L^2_tH^1}.
\]
Proposition 7 of Coron--Cr\'epeau \cite{CC} provides continuous linear maps
from $H\times L^1(0,T;H)$ to the mild solution $v\in X_T$ of
$v_t=Av+f$, $v(0)=h$, and to its left derivative in $L^2(0,T)$, with zero
right Neumann datum. The latter map defines $\mathcal Bv$ for rough data.
It agrees with the actual derivative for classical solutions. In particular,
\begin{equation}\label{eq:forced-estimate}
 \|v\|_{X_T}+\|\mathcal Bv\|_{L^2(0,T)}
 \le K_T\bigl(\|h\|_2+\|f\|_{L^1_tH}\bigr).
\end{equation}
Here and below $\|\cdot\|_2$ is the spatial $L^2$ norm.
The identities behind this estimate are
\begin{align}
 \|v(t)\|_2^2+\int_0^t|\mathcal Bv(s)|^2\,ds
 &=\|h\|_2^2+2\operatorname{Re}\int_0^t\langle f(s),v(s)\rangle\,ds,
 \label{eq:energy-linear}\\
 \frac{d}{dt}\int_0^L x|v|^2\,dx+3\|v_x\|_2^2
 &=\|v\|_2^2+2\operatorname{Re}\int_0^L xf\overline v\,dx.
 \label{eq:multiplier}
\end{align}
For smooth data these follow by integration by parts. The first gives
$\sup_t\|v(t)\|_2\le\|h\|_2+\|f\|_{L^1_tH}$ and the trace bound;
the integrated second identity gives the spatial derivative bound.
To pass to mild data, approximate $h$ in $H$ by $C^\infty_c(0,L)$ and
$f$ in $L^1_tH$ by finite sums of smooth time functions with spatial values
in $C^\infty_c(0,L)$. The continuous maps of \cite[Proposition 7]{CC} give
strong convergence in $C_tH$, $L^2_tH^1$ and $L^2$ for the left trace.
Every term in the integrated identities converges, uniformly in its upper
endpoint. Zero Dirichlet traces pass in $H^1$, so the limit belongs to $X_T$.
Thus \eqref{eq:energy-linear} holds for every $t\in[0,T]$, and
\eqref{eq:multiplier} holds in its integrated form and almost everywhere.
This also identifies the mild solution with
$S(t)h+\int_0^t S(t-s)f(s)\,ds$.
The hidden trace respects time shifts: prove
$\mathcal BS(s)S(t)h=\mathcal BS(s+t)h$ first for $h\in D(A)$ and then use
its continuous extension in $L^2(0,T)$. We shall use this identity for fixed
$t\ge0$, without assigning a boundary derivative to an arbitrary state in $H$.

We also need graph-domain control of this observation. If $f=-Av=v'''+v'$,
then, with $a=v'(0)$ and $b=v''(0)$,
\[
 v(x)=a\sin x+b(1-\cos x)+F(x),\qquad
 F(x)=\int_0^x(1-\cos(x-s))f(s)\,ds.
\]
This formula has the prescribed three initial derivatives and satisfies the
ODE. Uniqueness with zero initial derivatives follows by differentiating
$|v'|^2+|v''|^2$ for the homogeneous equation.
The convolution kernels bound $F,F',F''$ in $L^2$, and $F'''=f-F'$.
The Gram matrix of $\sin x$ and $1-\cos x$ is invertible, hence their
coefficients are bounded by $C_L(\|v\|_2+\|F\|_2)$. Consequently
\begin{equation}\label{eq:graph}
 \|v\|_{H^3}+|\mathcal Bv|\le C_L(\|Av\|_2+\|v\|_2),\qquad v\in D(A).
\end{equation}
All endpoint values here use the continuous one-dimensional Sobolev
representatives; the corresponding Sobolev facts may be found in
\cite[Theorems 8.2, 8.8 and 8.12]{Brezis}.

\subsection{The three undamped modes}
The following explicit spectral description will also identify the subspace
on which linear decay fails.
\begin{lemma}\label{lem:modes}
The only imaginary eigenvalues of $A$ are $0,\pm ip$, where
\begin{equation}\label{eq:modes}
 p=\frac{120}{343},\qquad \phi(x)=1-\cos x,\qquad
 \psi(x)=2e^{-8ix/7}-13e^{3ix/7}+11e^{5ix/7}.
\end{equation}
Their eigenspaces span $M=\operatorname{span}_{\mathbb C}\{\phi,\psi,\overline\psi\}$.
They are mutually orthogonal and have zero derivative at both endpoints.
Every real trajectory in $M$ is
\begin{equation}\label{eq:U}
 U(t)=c\phi+ze^{-ipt}\psi+\overline z e^{ipt}\overline\psi,
 \qquad c\in\mathbb R,\quad z\in\mathbb C,
\end{equation}
with
\begin{equation}\label{eq:center-norm}
 \|U(0)\|_2^2=\frac{3L}{2}c^2+588L|z|^2.
\end{equation}
\end{lemma}
\begin{proof}
If $Av=-i\omega v$, $\omega\in\mathbb R$, then
\eqref{eq:dissipative} also gives $v'(0)=0$.
The spatial roots satisfy $\lambda^3+\lambda-i\omega=0$.
For distinct roots, the coefficients of $v=\sum a_j e^{\lambda_jx}$ and
of its endpoint values both lie in the one-dimensional nullspace of the
rows $(1,1,1)$ and $(\lambda_1,\lambda_2,\lambda_3)$.
Each coordinate of its null vector is nonzero. Thus
$e^{\lambda_1L}=e^{\lambda_2L}=e^{\lambda_3L}$.
Their real parts must be equal, and their sum is zero, so all roots are purely
imaginary. At a repeated root the roots are $ir,ir,-2ir$, with
$r=\pm1/\sqrt3$. The zero initial value and derivative give
$v=a e^{-2irx}+(-a+3ira x)e^{irx}$.
If $a\ne0$, its endpoint value requires
$e^{-3irL}=1-3irL$, impossible by absolute values. The remaining solution is zero.
These exponential descriptions exhaust the ODE solutions: the initial-jet
matrix is a Vandermonde matrix (or its double-root version), and a solution
with all initial jets zero vanishes by the differential inequality
$\frac d{dx}|(v,v',v'')|^2\le(1+|\omega|)|(v,v',v'')|^2$.

Write the three distinct roots as $iq_1,iq_2,iq_3$ in increasing order.
The common endpoint phase implies
$q_2-q_1=2\pi k/L$, $q_3-q_2=2\pi\ell/L$ for positive integers $k,\ell$.
Vi\`ete's relations give
$L=2\pi\sqrt{(k^2+k\ell+\ell^2)/3}$.
At $L=14\pi$ this is $k^2+k\ell+\ell^2=147$.
Here $k\le12$, and the discriminants for $k=1,\ldots,12$ are
\[
585,576,561,540,513,480,441,396,345,288,225,156.
\]
Only $576=24^2$, $441=21^2$ and $225=15^2$ are squares, giving
$(k,\ell)=(2,11),(7,7),(11,2)$. Each eigenspace has dimension one by the
same endpoint nullspace. Differentiating \eqref{eq:modes} gives
$A\phi=0$ and $A\psi=-ip\psi$, with all four endpoint values zero.
Finally every nonconstant exponential of frequency an integer multiple of
$1/7$ integrates to zero over $(0,14\pi)$. This proves the orthogonality and
\eqref{eq:center-norm}, since $2^2+13^2+11^2=294$.
\end{proof}

\subsection{Decay on the orthogonal complement}
Let $P$ be the orthogonal projection onto $M$ and $Q=I-P$.
\begin{lemma}\label{lem:stable}
The spaces $M$ and $M^\perp$ reduce $S(t)$. There are $K_s\ge1$ and
$\gamma>0$, depending only on $L$, such that
\begin{equation}\label{eq:stable}
 \|S(t)Q\|_{H\to H}\le K_s e^{-\gamma t}\qquad(t\ge0).
\end{equation}
Moreover $Q$ preserves $D(A)$ and $AQ=QA$ there.
\end{lemma}
\begin{proof}
For $T>0$ put $N_T=\{h\in H:\mathcal BS(\cdot)h=0\text{ in }L^2(0,T)\}$.
First, bounded initial states give a relatively compact family
$S(\cdot)h$ in $L^2(0,T;H)$. Indeed \eqref{eq:forced-estimate} bounds
these trajectories in $L^2_tH^1_0$.
For each $a>0$ compact spatial embedding and smooth density give a finite-dimensional
subspace $E_a\subset C^\infty_c(0,L)$ with
$\|u-P_{E_a}u\|_2\le a\|u\|_{H^1}$ for every $u\in H^1_0$.
For a smooth basis test $\varphi$ of this space the coefficient
$b(t)=\langle S(t)h,\varphi\rangle$ satisfies
$b'=-\langle (S(t)h)_x,\varphi+\varphi''\rangle$ in distributions.
It is bounded in $H^1(0,T)$ and hence compact in $C([0,T])$.
The finitely many coefficients are therefore compact, while the spatial
remainder is at most $a\|S(\cdot)h\|_{L^2_tH^1}$.
A diagonal subsequence as $a\downarrow0$ proves the assertion.

The kernel $N_T$ is closed by trace continuity. On this kernel
\eqref{eq:energy-linear} makes $h\mapsto S(\cdot)h$ an isometry up to
$\sqrt T$. Thus its unit ball is compact and $\dim N_T<\infty$.
The nonincreasing integer sequence $\dim N_n$ eventually stabilizes.
The nested subspaces $N_n\subset N_1$ therefore stabilize, so for some integer $T_0$,
\[
 N_{T_0}=\bigcap_{T>0}N_T=:E.
\]
Trace shifts give $S(t)E\subset E$. On this finite-dimensional space $S(t)$ is an
isometry, hence onto. We must also know that $E\subset D(A)$.
The averaged operators $T_h=h^{-1}\int_0^h S(s)|_E\,ds$ converge to $I$
in operator norm, because strong continuity is uniform in finite dimension.
They are invertible for small $h$, and the semigroup difference quotient gives
$T_hv\in D(A)$ and $AT_hv=(S(h)v-v)/h$.
Surjectivity yields the desired domain inclusion.
The real generator $D=A|_{E\cap H_{\mathbb R}}$ is skew-adjoint by
polarization of the norm identity. The self-adjoint nonnegative matrix
$-D^2$ has an orthonormal eigenbasis. On a positive eigenspace with eigenvalue
$\lambda$, a unit vector $u$ and $Du/\sqrt\lambda$ form a rotation plane;
its orthogonal complement is $D$-invariant. The zero eigenspace is the kernel
of $D$. Thus the complexification of $E$ has imaginary eigenvectors of $A$,
and Lemma~\ref{lem:modes} implies $E\subset M$. The reverse inclusion
follows from the zero observations of the explicit modes, so $E=M$.

There is consequently a $\theta>0$ such that
\begin{equation}\label{eq:observability}
 \|\mathcal BS(\cdot)h\|_{L^2(0,T_0)}\ge\theta\|h\|_2
 \qquad(h\in M^\perp).
\end{equation}
Otherwise take unit $h_n\in M^\perp$ whose observations tend to zero.
By the compactness just proved their trajectories have a Cauchy subsequence.
Applying the time-averaged energy identity to $h_n-h_m$ gives
\[
 \|h_n-h_m\|_2^2\le T_0^{-1}\|S(\cdot)(h_n-h_m)\|_{L^2_tH}^2
 +\|\mathcal BS(\cdot)(h_n-h_m)\|_{L^2(0,T_0)}^2.
\]
The initial states converge to a unit vector in $M^\perp\cap N_{T_0}$,
which is impossible.
Because $S(t)$ is contractive and unitary on $M$, applying contractivity to
$m+ah$ for every $a\in\mathbb C$, $m\in M$, $h\in M^\perp$, forces
$\langle S(t)m,S(t)h\rangle=0$. This proves the reducing property.
Decrease $\theta$ if necessary to $0<\theta\le1/2$. Energy and
\eqref{eq:observability} give $\|S(T_0)Q\|\le q:=\sqrt{1-\theta^2}<1$.
Iteration and contractivity between these times give \eqref{eq:stable},
with $K_s=q^{-1}$ and $\gamma=-\log q/T_0$.
Commutation of $P,Q$ with $S(t)$ and the generator difference quotient prove
the last assertion.
\end{proof}

\section{A coercive quadratic boundary observation}\label{sec:quadratic}
The linear trajectories in $M$ lose no energy. Their quadratic forcing must
therefore produce the first observable loss. We construct a response whose
boundary derivative is periodic, while allowing its unobserved part to grow.

\subsection{Separating the stable response from the center motion}
Recall the real trajectory $U(t)=c\phi+ze^{-ipt}\psi+\overline z e^{ipt}\overline\psi$ from \eqref{eq:U}, with $c\in\mathbb R$ and $z\in\mathbb C$. Put
\begin{equation}\label{eq:force}
 F=-UU_x=\sum_{j=-2}^2e^{-ijpt}F_j.
\end{equation}
The coefficients are smooth functions, quadratic in $c,z,\overline z$.
For $j=-2,\ldots,2$ define
\begin{equation}\label{eq:resolvent}
 R_j f=\int_0^\infty e^{ijpt}S(t)Qf\,dt.
\end{equation}
The stable estimate makes this a norm-convergent integral in $QH$, bounded
by $K_s\gamma^{-1}\|Qf\|_2$. Its graph-domain meaning follows directly from
\[
 S(h)R_jf=e^{-ijph}\left(R_jf-\int_0^h e^{ijps}S(s)Qf\,ds\right).
\]
Taking the right difference quotient proves
\begin{equation}\label{eq:inverse}
 R_jf\in QD(A),\qquad (-ijp-A)R_jf=Qf.
\end{equation}
This inverse is unique on $QH$, since a nonzero homogeneous solution would
contradict \eqref{eq:stable}. Estimate \eqref{eq:graph} bounds its boundary
trace, and $R_{-j}(\overline f)=\overline{R_jf}$.

Set
\begin{equation}\label{eq:response}
 W_s(t)=\sum_{j=-2}^2e^{-ijpt}R_jF_j,\qquad
 W_c(t)=\int_0^tS(t-s)PF(s)\,ds,\qquad W=W_s+W_c.
\end{equation}
These graph-domain functions solve $W_t=AW+F$.
Although $W_c$ can have secular terms, it lies in $M$ and
$\mathcal BW_c=0$. Therefore
\begin{equation}\label{eq:G}
 G(t)=\mathcal BW(t)=\sum_{j=-2}^2e^{-ijpt}\mathcal BR_jF_j
\end{equation}
is periodic with period $P_0=2\pi/p$.
Moreover $W(0)=W_s(0)$ is bounded in $H$, uniformly for $\|U(0)\|_2\le1$.
Only this initial bound, rather than an all-time bound on $W_c$, will be needed.

\subsection{The constant and double-frequency coefficients}
The zero-mode self-interaction has a stationary response
\begin{equation}\label{eq:w00}
 w_{00}(x)=\tfrac12(x-L)\sin x+\tfrac1{12}(\cos2x-1).
\end{equation}
It belongs to $D(A)$ and satisfies
$w_{00}'''+w_{00}'=-\phi\phi'$, $\mathcal Bw_{00}=-L/2$.
Its forcing is orthogonal to $M$: integer frequencies do not meet those of
$\psi,\overline\psi$, and $\int\phi^2\phi'=0$.
Thus $Qw_{00}$ is its stable response and has the same boundary derivative.

The oscillatory mean forcing is $-|z|^2(|\psi|^2)'$, where
\begin{equation}\label{eq:psi-square}
 |\psi|^2=294-52\cos(11x/7)+44\cos(13x/7)-286\cos(2x/7).
\end{equation}
The differentiated frequencies $\pm2/7,\pm11/7,\pm13/7$ are disjoint from
those of every element of $M$, so this forcing is orthogonal to $M$.
Writing $|\psi|^2=\sum d_qe^{iqx}$, the function
\[
 h(x)=\sum_{q\ne0}\frac{-d_q}{1-q^2}(e^{iqx}-1)
\]
satisfies $h'''+h'=-(|\psi|^2)'$ and all three domain boundary conditions.
Indeed $qL\in2\pi\mathbb Z$, and $d_q=d_{-q}\in\mathbb R$ makes its endpoint
derivative sum vanish. Hence $\mathcal Bh=0$, and projection onto $QH$
does not change this observation. The constant term of $G$ is exactly $-Lc^2/2$.

At frequency $2p$, let $v_2$ be the solution of
\begin{equation}\label{eq:v2}
 v_2'''+v_2'-2ip v_2=-\psi\psi',\qquad
 v_2(0)=v_2(L)=v_2'(L)=0,
\end{equation}
and set $\beta_2=\mathcal Bv_2$.
Existence and uniqueness follow by combining \eqref{eq:inverse} on $M^\perp$
with the invertible center block, whose factors are $-2ip,-ip,-3ip$.
Its full solution and its stable component have the same observation.
Also put $\beta_1=\mathcal BR_1[-(\phi\psi)']$, a finite number by
\eqref{eq:graph}. Expansion of \eqref{eq:force} now gives
\begin{equation}\label{eq:G-coefficients}
 G(t)=-\tfrac L2c^2+cz\beta_1e^{-ipt}
 +c\overline z\,\overline{\beta_1}e^{ipt}
 +z^2\beta_2e^{-2ipt}+\overline z^{\,2}\overline{\beta_2}e^{2ipt}.
\end{equation}
We need no nonvanishing assertion for $\beta_1$.

\begin{lemma}\label{lem:beta}
The coefficient $\beta_2$ is nonzero.
\end{lemma}
\begin{proof}
Write $q=(-8/7,3/7,5/7)$ and $a=(2,-13,11)$.
Since $q_j-q_j^3=p$, a particular solution of \eqref{eq:v2} is
\[
 \chi(x)=\frac16\left(\sum_{j=1}^3\frac{a_j}{q_j}e^{iq_jx}\right)^2.
\]
For verification, if $\xi=q_j+q_k$ then
$\xi-\xi^3-2p=-3q_jq_k(q_j+q_k)$; applying the differential operator to the
ordered quadratic sum gives precisely $-\psi\psi'$.
At both endpoints $\chi=d$ and $\chi'=0$, with
\begin{equation}\label{eq:d}
 d=\frac16\left(-\frac74-\frac{91}3+\frac{77}5\right)^2
 =\frac{1002001}{21600}>0.
\end{equation}
Let $\rho>1$ solve $\rho^3-\rho=240/343$. Monotonicity and rational
substitution give
\begin{equation}\label{eq:rho}
 \frac{31}{25}<\rho<\frac54.
\end{equation}
The spatial roots of \eqref{eq:v2} are
$\lambda_0=-i\rho$ and $\lambda_\pm=\pm\sigma+i\rho/2$, where
$\sigma=\sqrt{3\rho^2/4-1}>0$. Put $E_j=e^{\lambda_jL}$; their product is one.
The homogeneous correction $v_2-\chi=\sum b_je^{\lambda_jx}$ satisfies
\[
 \mathsf Qb=(-d,-d,0)^T,\qquad
 \mathsf Q=\begin{pmatrix}
 1&1&1\\ E_0&E_+&E_-\\
 \lambda_0E_0&\lambda_+E_+&\lambda_-E_-
 \end{pmatrix}.
\]
This matrix is invertible by the uniqueness already established for
\eqref{eq:v2}. Expanding its adjugate and using $E_0E_+E_-=1$ gives
\begin{equation}\label{eq:beta-formula}
 \beta_2=-\frac d{\det\mathsf Q}
 \sum_{j=0}^2\lambda_j(\lambda_{j+2}-\lambda_{j+1})(E_j+E_j^{-1}),
\end{equation}
where indices are cyclic in the order $(0,+,-)$.
To examine the numerator, put $\vartheta=\rho L/2=7\pi\rho$,
$v=\sigma L$ and $K=3\rho^2/2-1$. The three cyclic prefactors in
\eqref{eq:beta-formula} are $2i\rho\sigma$, $K-i\rho\sigma$ and
$-K-i\rho\sigma$. Since
$E_0+E_0^{-1}=2\cos(2\vartheta)$ and
$E_\pm+E_\pm^{-1}=2\cosh(\pm v+i\vartheta)$, the sum is
\begin{equation}\label{eq:bracket}
 4i\bigl[\rho\sigma\cos(2\vartheta)
 +K\sinh(v)\sin\vartheta-\rho\sigma\cosh(v)\cos\vartheta\bigr].
\end{equation}
By \eqref{eq:rho}, $\vartheta$ modulo $2\pi$ is $\pi(7\rho-8)$, with
$17/25<7\rho-8<3/4$. Thus $\sin\vartheta>1/2$ and $\cos\vartheta<0$.
Furthermore $\sigma^2>383/2500$, so $v>2$;
$\sigma<1/2$ and $\rho\sigma<5/8$; and
$K>1633/1250>13/10$.
The real bracket in \eqref{eq:bracket} is therefore larger than
\[
 -\frac58+\frac{13}{10}\cdot1\cdot\frac12=\frac1{40}>0.
\]
Thus \eqref{eq:beta-formula} cannot vanish.
\end{proof}

\begin{proposition}\label{prop:coercivity}
There is an $\eta>0$ such that every real center trajectory with
$1/2\le\|U(0)\|_2\le1$ satisfies
\begin{equation}\label{eq:coercivity}
 \|G\|_{L^2(T,T+P_0)}\ge\eta\qquad\text{for every }T\ge0.
\end{equation}
\end{proposition}
\begin{proof}
Integrating \eqref{eq:G-coefficients} over a full period gives
\[
 \int_T^{T+P_0}|G(t)|^2\,dt
 =P_0\left(\frac{L^2}{4}c^4+2|\beta_1|^2c^2|z|^2
                 +2|\beta_2|^2|z|^4\right).
\]
The cross terms have distinct temporal frequencies, hence their integrals
vanish for every starting time $T$. By \eqref{eq:center-norm}, the lower
shell bound gives $c^2+|z|^2\ge1/(4\cdot588L)$. Since
$c^4+|z|^4\ge(c^2+|z|^2)^2/2$, one may choose
\[
 \eta^2=\frac{P_0}{32(588L)^2}\min\{L^2/4,2|\beta_2|^2\}>0.
\]
This one constant covers the zero mode, the oscillatory pair, and every
mixed direction and phase.
\end{proof}

\section{Uniform nonlinear energy loss and decay}\label{sec:nonlinear}
We now compare a small nonlinear solution with the quadratic responses on
one fixed time interval. All observation times will be chosen before the
smallness threshold and the initial datum.

Small-data global well-posedness in the usual energy class is given by
\cite[Lemma 3.2]{NX25} with zero Neumann control. The mild realization is the
one in \cite[Proposition 7 and Appendix A]{CC}: on each finite interval,
$y=S(\cdot)y_0-\int_0^{\cdot}S(\cdot-s)(yy_x)(s)\,ds$.
The source's usual energy solution has precisely this meaning; see also
\cite[Lemma 5.4]{CKN}. Only existence and uniqueness are used from the
nonlinear well-posedness statement.

For such a real solution $yy_x\in L^1_tH$ on every finite interval, by
$H^1\hookrightarrow L^\infty$ and time H\"older. We may apply the linear
identities with this exact force. Since
$\int_0^L y^2y_x\,dx=0$ for $y\in H^1_0$, we obtain
\begin{equation}\label{eq:energy-nonlinear}
 E(t)+\int_0^t|\mathcal By(s)|^2\,ds=E(0),\qquad E(t)=\|y(t)\|_2^2.
\end{equation}
The weighted identity becomes
\[
 \frac d{dt}\int_0^L xy^2\,dx+3\|y_x\|_2^2
 =\|y\|_2^2+\frac23\int_0^L y^3\,dx.
\]
These polynomial integration-by-parts identities hold for $H^1_0$ by its
smooth density and the continuous $H^1\to L^\infty$ embedding.
If $\varepsilon=\|y_0\|_2\le1$, energy and
$\left|\int y^3\right|\le\sqrt L\,\varepsilon^2\|y_x\|_2$, followed by
Young's inequality, give for each fixed $T$
\begin{equation}\label{eq:nonlinear-bound}
 \|y\|_{X_T}\le D_T\varepsilon.
\end{equation}
In particular this estimate requires no initial Sobolev regularity.
The bilinear estimate
\begin{equation}\label{eq:product}
 \|ab_x\|_{L^1(0,T;H)}
 \le C_L\|a\|_{L^2_tH^1}\|b\|_{L^2_tH^1}
\end{equation}
and \eqref{eq:forced-estimate} yield
\begin{equation}\label{eq:first-expansion}
 \|y-S(\cdot)y_0\|_{X_T}
 +\|\mathcal B(y-S(\cdot)y_0)\|_{L^2(0,T)}
 \le D'_T\varepsilon^2.
\end{equation}

\begin{proposition}\label{prop:loss}
There exist $T_*>0$, $\delta>0$ and $\kappa>0$, depending only on $L$,
such that every real solution with $\|y_0\|_2\le\delta$ satisfies
\begin{equation}\label{eq:loss}
 E(T_*)\le E(0)-\kappa E(0)^2.
\end{equation}
\end{proposition}
\begin{proof}
The zero datum is immediate from \eqref{eq:energy-nonlinear}; assume
$\varepsilon>0$. Choose $T_1$ so large that energy and \eqref{eq:stable} give
$\|\mathcal BS(\cdot)h\|_{L^2(0,T_1)}\ge\frac12\|h\|_2$ for $h\in M^\perp$.
Choose $\beta>4D'_{T_1}$. The scale $\beta\varepsilon^2$ separates initial
stable components already visible in the linear approximation from data whose
first relevant boundary signal is quadratic.

If $\|Qy_0\|_2\ge\beta\varepsilon^2$, then \eqref{eq:first-expansion} gives
\begin{equation}\label{eq:case-one}
 \|\mathcal By\|_{L^2(0,T_1)}
 \ge\tfrac12\|Qy_0\|_2-D'_{T_1}\varepsilon^2
 \ge\tfrac\beta4\varepsilon^2.
\end{equation}
Here $\mathcal BS(t)Py_0=0$.

Suppose instead that $\|Qy_0\|_2<\beta\varepsilon^2$, and put
\[
 U(t)=\varepsilon^{-1}S(t)Py_0,\qquad r=\varepsilon^{-2}Qy_0.
\]
For $\varepsilon\le1/(2\beta)$, orthogonality gives
$1/2\le\|U(0)\|_2\le1$ and $\|r\|_2\le\beta$.
Let $V_t=AV-UU_x$, $V(0)=r$. The response \eqref{eq:response} gives
\begin{equation}\label{eq:V}
 V(t)=W(t)+S(t)(r-W(0)).
\end{equation}
Both $r$ and $W(0)$ lie in $M^\perp$.
Let $A_0$ bound $\|W(0)\|_2$ on the shell above.
Choose a fixed $T_2$ such that
$K_s e^{-\gamma T_2}(\beta+A_0)\le\eta/2$.
The boundary trace of the homogeneous term in \eqref{eq:V} on
$(T_2,T_2+P_0)$ is at most its state norm at $T_2$, by the shifted energy
identity. Proposition~\ref{prop:coercivity} therefore gives
\begin{equation}\label{eq:V-trace}
 \|\mathcal BV\|_{L^2(T_2,T_2+P_0)}\ge\eta/2.
\end{equation}

Set $T_* =\max\{T_1,T_2+P_0\}$.
Equations \eqref{eq:first-expansion} and \eqref{eq:forced-estimate} imply
$\|y-\varepsilon U\|_{X_{T_*}}\le C\varepsilon^2$, uniformly on this case.
Indeed the omitted linear term is $S(t)Qy_0$ with initial norm at most
$\beta\varepsilon^2$. Using \eqref{eq:product} on
$(y-\varepsilon U)y_x+\varepsilon U(y-\varepsilon U)_x$ gives
\[
 \|yy_x-\varepsilon^2UU_x\|_{L^1(0,T_*;H)}\le C\varepsilon^3.
\]
The function $y-\varepsilon U-\varepsilon^2V$ has zero initial value and this
force with the opposite sign. Estimate \eqref{eq:forced-estimate} thus bounds
its boundary trace by $D\varepsilon^3$, for one $D\ge1$ fixed in advance.
Since $\mathcal BU=0$, \eqref{eq:V-trace} implies, if
$\varepsilon\le\eta/(4D)$,
\begin{equation}\label{eq:case-two}
 \|\mathcal By\|_{L^2(0,T_*)}\ge\eta\varepsilon^2/4.
\end{equation}
The possible secular growth of $W_c$ causes no additional estimate: this
comparison uses a fixed interval, and its observation vanishes exactly.

Choose $\kappa=\min\{\beta^2/16,\eta^2/16\}>0$.
Choose $\delta$ smaller than the global well-posedness threshold, $1$,
$1/(2\beta)$, $\eta/(4D)$ and $1/\sqrt{2\kappa}$.
Every constant and time has now been fixed independently of the initial
datum. Squaring \eqref{eq:case-one} or \eqref{eq:case-two} and using
\eqref{eq:energy-nonlinear} proves \eqref{eq:loss}.
\end{proof}

\begin{proof}[Proof of Theorem~\ref{thm:main}]
Energy keeps every shifted trajectory inside the same small ball, and the
equation is autonomous. Thus $E_n=E(nT_*)$ satisfies
$E_{n+1}\le E_n(1-\kappa E_n)$.
If some $E_n=0$, the solution subsequently vanishes by
\eqref{eq:energy-nonlinear}. Otherwise, since $\kappa\delta^2\le1/2$,
\[
 \frac1{E_{n+1}}\ge\frac1{E_n}+\kappa,\qquad
 E_n\le\frac{E_0}{1+\kappa nE_0}\le\frac1{\kappa n}\quad(n\ge1).
\]
For $t\ge2T_*$, $n=\lfloor t/T_*\rfloor\ge t/(2T_*)$ and monotonicity gives
$E(t)\le2T_*/(\kappa t)$. On $[0,2T_*]$ we have $E(t)\le\delta^2$.
Hence the theorem holds with
\[
 C=\max\{\delta\sqrt{1+2T_*},\sqrt{(2T_*+1)/\kappa}\}.
\]
This includes zero data, arbitrarily small nonzero data, and every time
between the observation windows.
\end{proof}
\section*{AI-use statement}
Substantial OpenAI Codex assistance was used in mathematical exploration,
calculations, source checking and preparation of this manuscript. Two additional
fresh-context Codex agents supplied reader and referee reports, followed by an
author revision. All three agents used the same vendor and shared host filesystem.
Independent mathematical review is requested.

\begin{thebibliography}{9}
\bibitem{Brezis}
H. Brezis, \emph{Functional Analysis, Sobolev Spaces and Partial Differential
Equations}, Springer, 2011. \href{https://doi.org/10.1007/978-0-387-70914-7}{doi:10.1007/978-0-387-70914-7}.
\bibitem{CC}
J.-M. Coron and E. Cr\'epeau, Exact boundary controllability of a nonlinear
KdV equation with critical lengths, \emph{J. Eur. Math. Soc.} \textbf{6}
(2004), 367--398. \href{https://doi.org/10.4171/JEMS/13}{doi:10.4171/JEMS/13}.
\bibitem{CKN}
J.-M. Coron, A. Koenig and H.-M. Nguyen, On the small-time local
controllability of a KdV system for critical lengths,
\emph{J. Eur. Math. Soc.} \textbf{26} (2024), no. 4, 1193--1253. \href{https://doi.org/10.4171/JEMS/1307}{doi:10.4171/JEMS/1307}.
Numbered source locations refer to the \href{https://ems.press/content/serial-article-files/31384}{EMS online-first version}.
\bibitem{Nguyen}
H.-M. Nguyen, Decay for the nonlinear KdV equations at critical lengths,
\emph{J. Differential Equations} \textbf{295} (2021), 249--291.
\href{https://doi.org/10.1016/j.jde.2021.05.057}{doi:10.1016/j.jde.2021.05.057}.
Numbered source locations refer to
\href{https://arxiv.org/abs/2012.08792v1}{arXiv:2012.08792v1} (2020).
\bibitem{NX25}
J. Niu and S. Xiang, Small-time local controllability of a KdV system for all
critical lengths, \href{https://arxiv.org/abs/2501.13640v2}{arXiv:2501.13640v2}
(2025), version of 16 December 2025.
\bibitem{NX26}
J. Niu and S. Xiang, Stability of a Korteweg--de Vries equation close to critical
lengths, \href{https://arxiv.org/abs/2603.27605v1}{arXiv:2603.27605v1} (2026).
\bibitem{Rosier}
L. Rosier, Exact boundary controllability for the Korteweg--de Vries equation
on a bounded domain, \emph{ESAIM Control Optim. Calc. Var.} \textbf{2}
(1997), 33--55. \href{https://doi.org/10.1051/cocv:1997102}{doi:10.1051/cocv:1997102}.
\end{thebibliography}
\end{document}
